%
% Einführung in die Mustererkennung - WS2013
% Abgabeprotokoll Exercise 1
%
%

%{{{ misc
\documentclass[a4paper,psfig,subfigure,epsfig,fleqn,amssmb,float,caption,fontenc,ausarbeitung]{article}

\usepackage{graphicx}
\usepackage{caption}
\usepackage{subcaption}

%Zitieren:
\usepackage[english]{babel}
%\usepackage[german]{babel}
\usepackage{babelbib} % für das Erstellen des Bibtex-Literaturverzeichnisses
\usepackage{cite}
%\selectbiblanguage{english}
%\selectbiblanguage{german}


\usepackage{url}

\usepackage[]{mcode}

\usepackage[pdftitle={Einfuehrung in die Mustererkennung, Exercise 1},
            pdfauthor={David Pfahler},
						pdfauthor={Matthias Gusenbauer},
						pdfauthor={Matthias Vigele},
            pdfsubject={Mustererkennung},
            pdfborder={0 0 0}]{hyperref}


%%%%%%%%%%%%%%%%%%%%%%%%%%%%%%%%
% Titlepage

\pagestyle{empty}


%set dimensions of columns, gap between columns, and paragraph indent

\setlength{\textheight}{24.7 cm}
\setlength{\columnsep}{1 cm}
\setlength{\textwidth}{16 cm}
%\setlength{\footheight}{0.0 cm}
\setlength{\topmargin}{0.0 cm}
\setlength{\headheight}{0.0 cm}
\setlength{\headsep}{-0.3 cm}
\setlength{\oddsidemargin}{0.0 cm}
\setlength{\parindent}{0.7 cm}
\setlength{\mathindent}{0mm}

% set page counter if document is part of proceedings
\setcounter{page}{1}
\renewcommand{\floatpagefraction}{0.9}
\renewcommand{\textfraction}{0.1}

%\renewcommand{\captionlabelfont}{\fontfamily{phv}\fontseries{bx}\fontsize{10}{10pt}\selectfont}
%\renewcommand{\captionfont}{\fontfamily{phv}\fontsize{10}{12pt}\selectfont}
%\setlength{\captionmargin}{0.5 cm}

\makeatletter
\makeatother
\def\RR{\hbox{I\kern-.2em\hbox{R}}}


\begin{document}

%don't want date printed
\date{}

%make title bold and 14 pt font (Latex default is non-bold, 16pt) 
\title{~\\
  ~\\
  \fontsize{14}{14pt} \bf Abgabedokument Exercise 1
	 ~\\
  \fontsize{12}{12pt} \bf Einführung in die Mustererkennung 186.840 WS 2013}

%for single author 
\author{~\\
  ~\\
  \fontsize{12}{12pt}
  {\bf David Pfahler, Matthias Gusenbauer, Matthias Vigele}\\
  1126287, 1125577, 1126171
  ~\\ ~\\ ~\\
  \normalsize
  {\bf ABSTRACT} \\ 
  \noindent
  \hspace{0.2cm}
  \begin{minipage}[c]{15cm}
  \normalsize This is my abstract.  This is my abstract.  This is my
    abstract.  This is my abstract.  This is my abstract.  This is my
    abstract.  This is my abstract.  This is my abstract.  This is my
    abstract.  This is my abstract.  This is my abstract.  This is my
    abstract.\\
  \end{minipage}
  ~\\ ~\\ ~\\
  \normalsize
  {\bf KEYWORDS} \\ 
  \normalsize
  Keyword1, Keyword2, Keyword3, Keyword4, Keyword5, Keyword6.
  }

\maketitle

%I don't know why I have to reset thispagestyle, but otherwise get page numbers 
\normalfont
\thispagestyle{empty}

%%%%%%%%%%%%%%%%%%%%%%%%%%%%%%%%%%%%%%%%%%%%%%%%%%%%%%%%%%%%%%%%%%%%%%%%%%%%%%%%
% CONTENT

\section{Feature Extraction}
\label{sec:featureExtraction}

Just an matlab example:

\begin{lstlisting}
a = 1;    % MATLAB comment
b = 2;
c = a^2 + b^2;
feature = regionprops(inputImage, propname);
\end{lstlisting}


\subsection{Input Images}
\label{sec:input}

\subsection{Features}
\label{sec:features}

In the following section the used features are described. To evaluate scale, 
transformation and rotation invariants a test image is used. Listing 
\ref{lst:calcRot} shows the used matlab function. Figure \ref{fig:heart} 
shows the test image. 

\begin{lstlisting}[caption=Calculate rotated image, label=lst:calcRot]
	rot = imresize(imrotate(image,45),0.5);
\end{lstlisting}
The calculation of the features of the two test images will differ, because 
of the lost of accuracy by the size reduction.

\subsubsection{Eccentricity}
\label{sec:eccentricity}

The Eccentricity is a scalar value that measures of how much the object 
deviates from a circle. It is calculated by the ratio of the distance between 
the foci of the ellipse and its major axis length. See Equation 
\ref{eq:eccentricity}
\begin{equation}
	\frac{2*\sqrt{\frac{D_{max}}{2}^2-\frac{D_{min}}{2}^2}}{D_{max}}
	\label{eq:eccentricity}
\end{equation}
$D_{max}$ is the major axis length and $D_{min}$ is the minor axis length of 
the best-fit ellipse.

\paragraph{Invariance:} The eccentricity feature is scale, transformation and 
rotation invariant. The Listing \ref{lst:eccentricity} provides an example. 
Figure \ref{fig:heart} shows the used binary images.

\begin{lstlisting}[caption=Calculate eccentricity in matlab, label=lst:eccentricity]
>> regionprops(image,'Eccentricity')

ans = 

    Eccentricity: 0.7073

>> regionprops(rot,'Eccentricity')

ans = 

    Eccentricity: 0.7080
\end{lstlisting}

\subsubsection{Roundness}
\label{sec:roundness}

Roundness is a scalar value that describe how the shape of an object 
approaches that of a circle. A value of 1 means that the object is a circle 
and a value of 0 means the object is a line segment. See Equation \ref{eq:roundness}
\begin{equation}
  \frac{4A(R)}{\pi D^{2}_{max}}
	\label{eq:roundness}
\end{equation}
$A(R)$ is the area of the labeled binary image and $D_{max}$ is the major axis 
length of the best-fit ellipse.

\paragraph{Invariance:} The Roundness feature is scale, transformation and 
rotation invariant. The Listing \ref{lst:roundness} provides an example. 
Figure \ref{fig:heart} shows the used binary images.

\begin{lstlisting}[caption=Calculate roundness in matlab, label=lst:roundness]
	>> calculateRoundness(image)

	ans = 

			Roundness: 0.6333
	
	>> calculateRoundness(rot)

	ans = 

			Roundness: 0.6326
\end{lstlisting}

\begin{figure}
	\centering
	\begin{subfigure}[b]{0.4\textwidth}
		\includegraphics[width=\textwidth]{img/heart.jpg}
	\end{subfigure}
	~ 
	\begin{subfigure}[b]{0.4\textwidth}
		\includegraphics[width=\textwidth]{img/invariance_rotation.jpg}
	\end{subfigure}
	\caption{A rotated, translated and scaled image}
	\label{fig:heart}
\end{figure}


\subsubsection{Solidity}
\label{sec:solidity}

The solidity is a scalar value specifying the proportion of the pixels in the 
convex hull that are also in the region. See Equation \ref{eq:solidity}

\begin{equation}
\frac{A(R)}{A(Co(R))}
\label{eq:solidity}
\end{equation}
$A(R)$ is the area of the labeled binary image and $Co(R)$ is the convex hull of the binary image

\paragraph{Invariance:} The solidity feature is scale, transformation and 
rotation invariant. The Listing \ref{lst:solidity} provides an example. 
Figure \ref{fig:heart} shows the used binary images.
\begin{lstlisting}[caption=Calculate solidity in matlab, label=lst:solidity]
>> regionprops(image,'Solidity')

ans = 

    Solidity: 0.9145

>> regionprops(rot,'Solidity')

ans = 

    Solidity: 0.9113
\end{lstlisting}
\section{Classification}
\label{sec:classification}

\subsection{Own Implementation}
\label{sec:own}

\subsubsection{Methods}
\label{sec:ownMethods}

\subsubsection{Results}
\label{sec:ownResults}

\subsection{k-NN Classifier}
\label{sec:kNN}

\subsubsection{Methods}
\label{sec:kNNMethods}

\subsubsection{Results}
\label{sec:kNNResults}

\section{Paper Summary}
\label{sec:paperSummary}
\cite{beyer1999nearest}

% Biblography


\fontsize{9}{10pt}
\bibliographystyle{plain}
\bibliography{literatur}

\end{document}

% vim:foldmethod=marker
